\documentclass[11pt]{article}
\usepackage[margin=1in]{geometry}

\usepackage{booktabs,amsmath,amssymb,hyperref}
\usepackage[font=small]{caption}
\title{Does reallocating the perturbation budget or the training samples help the weak classes?\\
\large A reproduction-based study of per-class adaptive adversarial training on top of RAMP (CIFAR-10)}
\author{Matoke Brian\\\small Draft for supervisor review, 5 October 2026}
\date{}
\begin{document}
\maketitle
\begin{abstract}
We test whether feedback-driven, per-class scheduling of the training perturbation budget $\varepsilon$ improves
union robustness over $\{\ell_\infty,\ell_2,\ell_1\}$ on CIFAR-10 for a PreActResNet-18, using the official RAMP
training code as the baseline and AutoAttack's APGD-CE and APGD-T attacks at \emph{fixed} $\varepsilon$ for every evaluation. Two budget-neutral
controllers (more $\varepsilon$ for classes that are already robust, or the reverse) and a budget-neutral
class-sampling fine-tune do not improve union robustness or the weakest classes over their controls; the reversed
controller reduced union robustness in both seeds. The weakest classes (bird, cat, deer, dog) are also poorly fit on training images,
which is consistent with an optimisation- or capacity-limited explanation, but the experiments are low-powered
(2 seeds, about 100 test images per class) and do not separate that explanation from a data-limited one.
We list the experiments needed to do so.
\end{abstract}

\section{Motivation and scope}
The starting point is the author's MSc thesis on Adaptive Adversarial Training (AAT)~\cite{thesis} that scheduled the perturbation budget and
the attacked norm from per-class feedback. An audit of that thesis (\texttt{docs/thesis\_audit.md}) found
inconsistent numbers and, most importantly, evaluation at the \emph{adaptive} $\varepsilon$, which inflates robust
accuracy. This study therefore re-tests the idea under a clean protocol and does not reproduce the thesis numbers.
The research question is whether feedback scheduling improves union robustness compared with fixed-schedule
multi-norm adversarial training (RAMP~\cite{ramp}).

\section{Setup}
\textbf{Model and baseline.} PreActResNet-18 trained with the official RAMP code (pinned commit) for 80 epochs from
scratch ($\lambda=5$, gradient projection), 2 seeds. The baseline reproduces the thesis Table~7.1 baseline to within one
point on every metric.\\
\textbf{Evaluation.} the APGD-CE and APGD-T attacks of AutoAttack~\cite{aa} (not the full suite: no FAB-T or Square) at fixed $\varepsilon=8/255$ ($\ell_\infty$), $0.5$
($\ell_2$), $12$ ($\ell_1$) on the first 1000 CIFAR-10 test images; union robustness is the per-sample AND over the three
norms. An $\varepsilon$ grid beyond the training budgets (APGD-CE only) measures robustness to unseen threat sizes.
Results are taken from \texttt{eval\_autoattack.json}, never from the 200-image monitor in the training logs.\\
\textbf{Arms.}
\emph{aat\_rel}/\emph{aat\_flip}: RAMP plus a per-class training-$\varepsilon$ controller for the $\ell_\infty$ and $\ell_1$ training attacks (the $\ell_2$ component of RAMP is not modified), a budget-neutral variant of the thesis update (Eq.~4.19): the set-point $\tau$ is replaced by the class-mean training-attack accuracy, $\alpha=0.05$, updated once per epoch, $\varepsilon_c$ started at nominal, clipped to $[0.5,1.5]\times$nominal and projected so that the class mean equals nominal. aat\_rel gives more $\varepsilon$ to classes with higher training robust accuracy, aat\_flip the reverse (a fairness-style allocation~\cite{frl}).
\emph{cls\_ctrl}/\emph{cls\_fb}: the locked epoch-80 RAMP weights continued for 10 epochs at RAMP's final-phase constant learning rate (0.005), resuming optimizer and GP state, with a
with-replacement sampler drawing as many images as the dataset has; cls\_ctrl samples classes uniformly, cls\_fb with weights
$w_c=\mathrm{clip}(1+(\bar A-A_c)/\bar A,\,0.5,\,1.5)$ (mean 1, so the number of images per epoch is unchanged), where $A_c$ is
the per-class training-attack robust accuracy averaged over the $\ell_\infty$ and $\ell_1$ attacks, the clipped weights are re-projected to mean 1, and $w=1$ in the first fine-tune epoch.
\emph{Checkpoint averaging}: uniform average of the epoch 60/70/80 weights (batch-norm statistics averaged, not recalibrated); the average spans the learning-rate drop at epoch 70 (epoch 60 was trained at 0.05).

\section{Results}
All tables are 2-seed means in percent. Per-class cells hold about 100 images (binomial standard error about 5 points),
and two seeds support no significance claim, so every difference below is descriptive.

\begin{table}[h]\centering\caption{Fixed-$\varepsilon$ AutoAttack, first 1000 test images.}
\begin{tabular}{lrrrrrr}\toprule
Run & Clean & $\ell_\infty$ & $\ell_2$ & $\ell_1$ & Union & Worst class \\\midrule
RAMP & 81.30 & 45.70 & 66.45 & 48.40 & 44.40 & 12.78 \\
aat\_rel & 82.90 & 45.80 & 66.85 & 48.55 & 44.65 & 10.83 \\
aat\_flip & 78.15 & 41.40 & 64.15 & 45.40 & 40.15 & 7.22 \\
cls\_ctrl & 82.50 & 45.00 & 65.40 & 48.50 & 43.45 & 12.78 \\
cls\_fb & 82.65 & 45.70 & 65.80 & 48.60 & 44.55 & 13.33 \\
\bottomrule\end{tabular}
\end{table}

\paragraph{$\varepsilon$ reallocation.} aat\_rel did not raise union over RAMP (44.65 vs 44.40), gained 1.6 clean points,
and lowered worst-class union in both seeds (10.0 and 11.7 vs 13.3 and 12.2). aat\_flip lowered union (40.15); both of its
seeds (42.0, 38.3) are below both RAMP seeds (45.7, 43.1). Neither supports the hypothesis that class-feedback
$\varepsilon$ narrows the class gap: worst-class union fell in both arms.

\paragraph{Sample reallocation.} Feedback class sampling was within noise of the uniform control (union 44.55 vs 43.45;
worst class 13.33 vs 12.78; mean union over bird, cat, deer and dog 21.1 vs 20.6), and neither arm exceeded locked epoch 80
by more than 0.2 points (cls\_fb 44.55, cls\_ctrl 43.45 vs 44.40). The $\varepsilon$ grid differs by at most one point between
the arms. The binomial standard error alone is about 5 points per class, so per-class effects of several points cannot be resolved, and a small benefit cannot be excluded.
Oversampling repeats images, so it does not test whether more \emph{distinct} data would help.

\paragraph{Checkpoint averaging.} Averaging lowered worst-class union in both seeds ($13.3\to10.7$, $12.2\to9.7$; the worst
class moved from deer to cat, about three images), changed union by $-1.3$ and $+2.8$ and clean accuracy by $-0.8$ and
$-2.2$. This does not support the claim that the worst-class number is mostly checkpoint noise, but given the size of the
effect and the unrecalibrated batch-norm statistics it does not refute it either.

\begin{table}[h]\centering\caption{Per-class union robustness (\%), 2-seed means; $\approx$100 images per cell.}
\begin{tabular}{lrrrrr}\toprule
Class & RAMP & aat\_rel & aat\_flip & cls\_ctrl & cls\_fb \\\midrule
plane & 50.0 & 49.5 & 52.9 & 46.6 & 46.1 \\
car & 59.0 & 63.5 & 52.8 & 59.6 & 57.9 \\
bird & 29.5 & 27.0 & 26.5 & 29.0 & 29.5 \\
cat & 13.1 & 12.1 & 9.2 & 16.0 & 16.5 \\
deer & 12.8 & 11.7 & 9.4 & 12.8 & 13.3 \\
dog & 27.9 & 23.8 & 36.6 & 24.4 & 25.0 \\
frog & 54.9 & 54.5 & 52.7 & 50.4 & 57.1 \\
horse & 63.2 & 63.7 & 57.8 & 62.3 & 63.2 \\
ship & 60.4 & 67.5 & 42.9 & 63.7 & 61.8 \\
truck & 65.6 & 65.1 & 55.5 & 62.4 & 66.5 \\
\bottomrule\end{tabular}
\end{table}

\paragraph{Class by norm.} For RAMP, in the 2-seed mean, union is at most 2.5 points below $\min(\ell_\infty,\ell_1)$ in every
class (per seed up to 4.4, deer seed 1) and 1.3 points overall. The weak classes are weak under every norm. For RAMP's weights,
better alignment across norms could add at most about 1.3 points of union; this does not bound what a different schedule could reach.

\paragraph{Unseen threat sizes.} Table~\ref{tab:unseen} shows the grid beyond the training budgets. aat\_rel is one to two points above
RAMP at five of the six budgets ($\ell_2$ at 1.0 is 0.2 below) and aat\_flip is below at every budget; only $\ell_\infty$ at 16/255 has non-overlapping seed ranges.
We regard this as a direction to test with a pre-specified protocol and more images, not as a finding.

\begin{table}[h]\centering\caption{Robust accuracy (\%) at budgets larger than training (APGD-CE), 2-seed means.}\label{tab:unseen}
\begin{tabular}{lrrrrrr}\toprule
Run & $\ell_\infty$ 12/255 & $\ell_\infty$ 16/255 & $\ell_2$ 1.0 & $\ell_2$ 1.5 & $\ell_1$ 18 & $\ell_1$ 24 \\\midrule
RAMP & 29.0 & 13.0 & 48.6 & 28.8 & 34.1 & 21.5 \\
aat\_rel & 31.0 & 14.5 & 48.4 & 29.8 & 35.3 & 22.9 \\
aat\_flip & 24.8 & 10.9 & 45.7 & 25.0 & 30.6 & 19.6 \\
cls\_ctrl & 29.2 & 13.0 & 49.0 & 28.2 & 34.2 & 21.5 \\
cls\_fb & 29.2 & 13.2 & 48.1 & 28.4 & 33.9 & 21.4 \\
\bottomrule\end{tabular}

\end{table}

\section{Why are the weak classes weak?}
To separate a data-limited from an optimisation- or capacity-limited explanation we evaluated the locked RAMP epoch-80 weights on the
first 1000 \emph{training} and test images (unaugmented; unweighted class means over 2 seeds). Union robustness of the weak four classes is
32.8 on training and 20.8 on test images (gap 12.0); for the other six classes it is 69.5 and 58.9 (gap 10.5). Clean accuracy of the weak
four on training images is 79.1, against 95.3 for the others. The weak classes are therefore poorly fit even on images the model was trained
on, and their train--test gap is not distinguishable from that of the other classes (cf.\ robust overfitting~\cite{rice}). Per class, cat (4--6
points) and dog (4--6) have small gaps while deer (20--28) and bird (11--17) have large ones, as do some strong classes; frog is 0 and $-7$.

\section{Interpretation and limits}
None of the three reallocation experiments raised union robustness or the weak classes. Reversing the $\varepsilon$ allocation lowered union (by about 4 points in both seeds) and worst-class union, but not the mean of the four weak classes (20.4 vs 20.8); aat\_rel's worst-class union was lower in both seeds (descriptive). The weak classes are also poorly fit on training images. This is consistent with an optimisation- or
capacity-limited explanation (D2), but equally with low power, narrow reallocation ranges, short fine-tuning and repeated images, and nothing here
changed model capacity or the number of distinct training images. D1 (the weak classes are data-limited) and D2 therefore remain open. D1 and D2
label this diagnosis only; they are not the hypotheses H1--H4 of the project README.\\
\textbf{Limits.} Two seeds; first 1000 test images (binomial standard error about 1.6 points overall, about 5 per class); one fine-tuning
recipe per method; one model size; one dataset; run-to-run and checkpoint-to-checkpoint union variation of up to about 4 points.

\section{Proposed next steps}
\begin{enumerate}
\item \textbf{Power.} Re-evaluate the final models on the full 10\,000-image test set (about ten times the 33-minute per-model cost, i.e.\ about 5.5 GPU-hours per model and about 55 for ten models, an estimate; or APGD-CE-only for the unseen grid), and add seeds (five or more) for RAMP, aat\_rel and the controls.
\item \textbf{Unseen threats.} Pre-specify the aat\_rel versus RAMP comparison on the $\varepsilon$ grid and a held-out norm, using the extra images and seeds, before claiming a benefit.
\item \textbf{Capacity.} Fine-tune a pre-trained WRN-28-10 with the RAMP recipe. Rough estimate on a Kaggle T4, scaling the measured PreActResNet-18 epoch time (about 850\,s) by the roughly 9-fold compute ratio and assuming 3 fine-tune epochs: about 6 hours of training plus 4--5 hours of AutoAttack evaluation per seed on CIFAR-10 alone, and about 12--14 hours of training (the project handoff note says 14+) with the official 500K auxiliary set, which doubles the images per epoch. These are estimates, not measurements; memory at batch size 128 is untested. Caveats: the CIFAR-10-only variant is not the official RAMP WRN recipe (that script requires the 500K pseudo-labelled set, for which no verified source is available), and the RobustBench Gowal2020 checkpoint was pre-trained with extra data, so this experiment mixes capacity with data and cannot by itself separate D1 from D2.
\item \textbf{Other datasets.} CIFAR-100 needs a new locked baseline and the full test set (100 images per class instead of 10 in the first 1000), so it is a generalisation check, not a cheap one.
\item \textbf{Full method.} The full thesis method (adaptive norm weights, hybrid loss, logit pairing) has not been run.
\end{enumerate}

\begin{thebibliography}{9}
\bibitem{thesis} B.~O.~Matoke. Adaptive Adversarial Training: A Curriculum-Based Approach to Robustness. MSc thesis, Department of Computer Science, University of Szeged, 2025.
\bibitem{ramp} E.~Jiang and G.~Singh. RAMP framework: Boosting adversarial robustness against multiple-$\ell_p$ perturbations. NeurIPS 2024.
\bibitem{aa} F.~Croce and M.~Hein. Reliable evaluation of adversarial robustness with an ensemble of diverse parameter-free attacks. ICML 2020.
\bibitem{frl} H.~Xu, X.~Liu, Y.~Li, A.~Jain and J.~Tang. To be robust or to be fair: towards fairness in adversarial training. ICML 2021.
\bibitem{rice} L.~Rice, E.~Wong and J.~Z.~Kolter. Overfitting in adversarially robust deep learning. ICML 2020.
\end{thebibliography}
\end{document}
